\section*{第 \ref{ch:b2:neurons-synapses} 章——神经元、动作电位与突触}

\begin{solution}{exo:b2:neurons-synapses:1}
\omterm{def:b2:neurons-synapses:neuron}{树突}接受\omterm{def:b2:neurons-synapses:synapse}{突触}输入；胞体含有细胞核，
并进行整合；\omterm{prop:b2:neurons-synapses:integration}{轴突起始段}决定是否发放；\omterm{def:b2:neurons-synapses:neuron}{轴突}传导冲动；
末梢释放递质。\omterm{def:b2:neurons-synapses:neuron}{髓鞘}是包裹\omterm{def:b2:neurons-synapses:neuron}{轴突}的胶质细胞膜，
几乎不含细胞质；它起绝缘作用，使
冲动在郎飞结之间跳跃式传播，速度快约一百倍，
而成本更低。
\end{solution}

\begin{solution}{exo:b2:neurons-synapses:2}
去极化达到阈值；上升相——电压门控钠通道
开放，形成正反馈；峰值接近 $E_{\mathrm{Na}}$——
钠通道失活；下降相——电压门控钾通道
开放；超射——钾通道仍开放，
电位接近 $E_{\mathrm K}$；随钾通道关闭而恢复。“全或无”是因为
正反馈要么进行到底，要么根本不发生；“单向”是因为
刚经过的膜处于\omterm{prop:b2:neurons-synapses:ap}{不应期}（钠通道已
失活）。
\end{solution}

\begin{solution}{exo:b2:neurons-synapses:3}
冲动到达末梢；电压门控钙通道开放；
钙离子触发囊泡融合；递质扩散穿过
\omterm{def:b2:neurons-synapses:synapse}{突触}间隙，与\omterm{def:b2:neurons-synapses:synapse}{突触}后受体结合；通道开放（或 G 蛋白
发生转换）；产生\omterm{def:b2:neurons-synapses:synapse}{突触后电位}；递质经酶降解或
经转运体回收；囊泡膜被回收。
\end{solution}

\begin{solution}{exo:b2:neurons-synapses:4}
电\omterm{def:b2:neurons-synapses:synapse}{突触}：无延迟，通常双向传递，无放大作用，无
抑制作用（电流只能通过）。化学\omterm{def:b2:neurons-synapses:synapse}{突触}：延迟约半毫秒，
单向传递，具有放大作用（一个冲动可释放数百个量子），
可产生抑制（氯离子或钾离子通道），并且
可塑。
\end{solution}

\begin{solution}{exo:b2:neurons-synapses:5}
由 $E = (61.5/z)\log_{10}(c_{\text{外}}/c_{\text{内}})$ 可得：
钾离子 \qty{-89}{mV}；钠离子 \qty{+61}{mV}；氯离子，$z =
-1$，$-61.5\log_{10} 11 = \qty{-64}{mV}$；钙离子，$z = 2$，
$30.75\times 4.3 = \qty{+132}{mV}$。在 \qty{-70}{mV} 时：
钾离子外流，钠离子内流，氯离子略微外流（膜
电位低于其\omterm{thm:b2:neurons-synapses:nernst}{平衡电位}），钙离子强烈内流。
\end{solution}

\begin{solution}{exo:b2:neurons-synapses:6}
$20 : 1$：$(20\times(-89) + 61)/21 = \qty{-82}{mV}$，静息态，
接近 $E_{\mathrm K}$。$1 : 20$：$(-89 + 20\times 61)/21 =
\qty{+54}{mV}$，峰值，接近 $E_{\mathrm{Na}}$。$50 : 1$：$(50\times
(-89) + 61)/51 = \qty{-86}{mV}$，超射，此时额外的钾
通道开放。
\end{solution}

\begin{solution}{exo:b2:neurons-synapses:7}
$\sqrt{500} = 22$：\qty{22}{m/s}。要达到 \qty{100}{m/s}，
无髓\omterm{def:b2:neurons-synapses:neuron}{轴突}需要直径 $d = 10^{4}$ \unit{\micro m}，
即一厘米。在有髓纤维中，传导一米需要 \qty{10}{ms}，
在鱿鱼巨\omterm{def:b2:neurons-synapses:neuron}{轴突}中需要 \qty{45}{ms}，在
\qty{1}{\micro m} 的纤维中则需整整一秒。
\end{solution}

\begin{solution}{exo:b2:neurons-synapses:8}
$m = \ln(200/74) = 1.0$; $P(1) = e^{-1} = 0.37$, $P(2) = 0.18$, $P(3)
= 0.06$。平均反应为 $1.0\times 0.5 = \qty{0.5}{mV}$。
\end{solution}

\begin{solution}{exo:b2:neurons-synapses:9}
$1/\qty{2}{ms} = 500$ 个冲动每秒。触觉的强度
由冲动的频率编码（以及被募集的\omterm{def:b2:neurons-synapses:neuron}{神经元}数目）；
冲动本身则完全相同。
\end{solution}

\begin{solution}{exo:b2:neurons-synapses:10}
$\lambda = \sqrt{R_m d/4R_i}$: \qty{0.5}{mm}, \qty{1.6}{mm},
\qty{11}{mm}。加\omterm{def:b2:neurons-synapses:neuron}{髓鞘}后（直径 \qty{10}{\micro m}）：$1.6\times\sqrt{200} =
\qty{22}{mm}$。一个长 \qty{1}{mm} 的节间仅损失 $1 - e^{-1/22}
= \qty{4}{\%}$ 的信号，因此下一个郎飞结很容易达到
阈值。跨越长 \qty{3}{mm} 的脱\omterm{def:b2:neurons-synapses:neuron}{髓鞘}段后，信号降至
$e^{-3/1.6} = \qty{15}{\%}$，即 \qty{100}{mV} 冲动的
\qty{15}{mV} 的阈值：传导失败或变得不可靠。
\end{solution}

\begin{solution}{exo:b2:neurons-synapses:11}
河豚毒素：无钠电流，无冲动，无传递。
四乙铵：无钾电流，冲动延长，每个冲动释放的
递质增多。无钙：冲动正常，
但无释放，无传递。筒箭毒碱：释放正常，受体
被阻断，无终板电位——导致麻痹。胆碱酯酶
抑制剂：乙酰胆碱持续存在，受体保持开放，肌肉
去极化后无法复极化——过量引起麻痹。
肉毒毒素：囊泡无法融合，无递质释放——导致麻痹。
\end{solution}

\begin{solution}{exo:b2:neurons-synapses:12}
延迟换来的是：单向传递；放大作用
（单个冲动可释放数百个量子，并能驱动
更大的细胞）；抑制作用，这是电\omterm{def:b2:neurons-synapses:synapse}{突触}做不到的；
以及可塑性，因为量子数和受体数可随使用而改变，
这正是回路学习的方式。电\omterm{def:b2:neurons-synapses:synapse}{突触}用于
同步与速度比运算更重要的场合：逃逸
回路、心肌以及某些抑制性网络。
\end{solution}

\begin{solution}{pb:b2:neurons-synapses:1}
\textbf{1.} $E_{\mathrm K} = 61.5\log_{10}(4/140) = \qty{-95}{mV}$;
$E_{\mathrm{Na}} = 61.5\log_{10}(145/12) = \qty{+67}{mV}$.
\textbf{2.} $(25\times(-95) + 67)/26 = \qty{-89}{mV}$.
\textbf{3.} $E_{\mathrm K} = 61.5\log_{10}(8/140) = \qty{-76}{mV}$;
$V = (25\times(-76) + 67)/26 = \qty{-71}{mV}$：起初更接近阈值，
因而兴奋性更强——但若持续，兴奋性会下降，原因在于
钠通道在去极化水平下失活。
\textbf{4.} 梯度会在数小时内逐渐消失，钠离子内漏，
钾离子外漏；电位向零漂移；若泵停止，
细胞内不能透过膜的阴离子会吸引氯离子和水进入，
细胞随之肿胀。
\textbf{5.} 每毫米的表面积为 $\pi\times 15\times 10^{-6}\times
10^{-3} = \qty{4.7e-8}{m^2}$; $C = \qty{4.7e-10}{F}$; $Q = CV =
\qty{4.7e-11}{C}$：$2.9\times 10^{8}$ 个离子。
\textbf{6.} 每毫米的体积为 $\pi(7.5\times 10^{-6})^{2}\times
10^{-3} = \qty{1.8e-13}{m^3} = \qty{1.8e-10}{L}$，其中含有 $12\times
10^{-3}\times 1.8\times 10^{-10} = 2.1\times 10^{-12}$ 摩尔，$1.3\times
10^{12}$ 个离子：变化量为 $2\times 10^{-4}$。
\textbf{7.} $2.9\times 10^{8}/3 \approx 10^{8}$ 个 ATP 每毫米
每个冲动。
\textbf{8.} $0.9/90 = \qty{10}{ms}$，每个方向相同。
\textbf{9.} 无髓：$\sqrt{1\times 15\times 10^{-6}/4} = \qty{1.9}{mm}$；
有髓：$\times\sqrt{250} = \qty{31}{mm}$。
\textbf{10.} 无髓：$e^{-1.5/1.9} = 0.45$，\qty{45}{mV}；有髓：
$e^{-1.5/31} = 0.95$，\qty{95}{mV}。两者均超过 \qty{15}{mV}
的阈值；但无髓\omterm{def:b2:neurons-synapses:neuron}{轴突}沿途必须给整个膜充电，
这正是其缓慢的原因；而在\omterm{def:b2:neurons-synapses:neuron}{髓鞘}之下，几乎
没有信号损失，也几乎无需充电。
\textbf{11.} $\sqrt{15} = \qty{3.9}{m/s}$：单程 \qty{0.23}{s}，
往返几乎需要半秒——对于必须应对绊倒的反射来说，
太慢了。
\textbf{12.} $e^{-4/1.9} = 0.12$：远端郎飞结处为 \qty{12}{mV}，低于
阈值——冲动被阻断，反射丧失。
\textbf{13.} 冲动之后，钠通道失活
约一到两毫秒，因此刚经过的膜不能
再次被兴奋，波只能向前传播。
\textbf{14.} $m = \ln(300/111) = 1.0$.
\textbf{15.} $P(1) = 0.37$，$P(2) = 0.18$，$P(3) = 0.06$：在 300 次试验中，约 110、55
和 18 次。
\textbf{16.} $1.0\times 0.4 = \qty{0.4}{mV}$.
\textbf{17.} $P(0) = e^{-60} \approx 10^{-26}$：几乎不可能；平均反应
为 \qty{24}{mV}，高于 \qty{15}{mV} 的阈值：单个此类\omterm{def:b2:neurons-synapses:synapse}{突触}
即可使运动\omterm{def:b2:neurons-synapses:neuron}{神经元}发放。
\textbf{18.} 二十个同时活动的\omterm{def:b2:neurons-synapses:synapse}{突触}叠加各自的电流
（空间\omterm{prop:b2:neurons-synapses:integration}{总和}），随着电位接近
通道的反转电位，呈次线性增长；即使每个\omterm{def:b2:neurons-synapses:synapse}{突触}的 $m = 1$，它们也能
产生 \qty{8}{mV}，而在正常 $m$ 值下远超阈值。
\textbf{19.} 在 $E_{\mathrm{Cl}} = V_{\text{静息}}$ 时开放氯离子通道，
增加了电导而不改变电位；兴奋性
电流此时流经漏电更多的膜，产生较小的
去极化——即分流抑制。
\textbf{20.} $10 + 0.6 + 10 + 0.6 + 5 = \qty{26}{ms}$，与
实测值 \qty{30}{ms} 相比。
\textbf{21.} 肌梭中感受器本身的激活、
\omterm{def:b2:cell-differentiation:myogenesis}{肌纤维}\omterm{prop:b2:neurons-synapses:ap}{动作电位}沿纤维的传播
（以米每秒的速度传播，跨越厘米量级的距离），以及
力出现之前的机械潜伏期。
\textbf{22.} 每个\omterm{def:b2:neurons-synapses:synapse}{突触}增加半毫秒，每个中间\omterm{def:b2:neurons-synapses:neuron}{神经元}
增加一个细胞；但中间\omterm{def:b2:neurons-synapses:neuron}{神经元}使信号得以反转
（抑制拮抗肌），可与其他信号整合，并受
大脑调制——使这一反射可以被塑造。
\textbf{23.} $90 - 6.2 = \qty{84}{ms}$ 用于传导 \qty{1.8}{m}：
约 \qty{21}{m/s}。
\textbf{24.} 在各电压下，内向电流均减半：阈值
升高，冲动变小、变慢，每个
郎飞结处的安全系数下降，冲动可能失败；反射减弱
或消失。
\textbf{25.} \omterm{thm:b2:neurons-synapses:chord}{静息电位} \qty{-89}{mV}；每个方向 \qty{10}{ms}；
$m = 1.0$；计算得出的潜伏期 \qty{26}{ms}。
\end{solution}
